\begin{afterword}
\summaryhead

The post is about reasons found after a conclusion has been reached. Its first example is Christians who no longer really believe but still revere the Bible, now as a source of ethics or as literature; the author replies that \textsc{The Lord of the Rings} and even \textsc{Harry Potter} are better on both counts. Its second is a friend who defends a \$1 million laptop by saying it stimulates the economy, when 5,000 cheap laptops for children would do more for the same goal.

Drawing on ``The Bottom Line,'' the author argues that only the real causes of a decision determine how good it is. A justification added later changes nothing unless it is a real search that could have changed the conclusion, and people overestimate how likely they are to change their minds. If an indefensible process produced a choice, it would be a suspicious coincidence for that choice to be the best answer under a respectable criterion. The post ends: ``It's improbable that you used mistaken reasoning, yet made no mistakes.''

\responsehead

The post's central argument is sound. A reason found after a decision says nothing about how the decision was made, and if the process that made it had nothing to do with the criterion now invoked, it would be a coincidence for the choice to be the criterion's best answer. The laptop example shows both the error and a way to test for it: look for an option that does better under the stated goal. The claim that people overestimate their readiness to change their minds is unsourced in this post, but the author had given evidence for it a month earlier. What the post describes is rationalizing in the ordinary sense of the word, the process the author had objected to naming that way in ``Rationalization'' (see the notes there).

The Bible example is weaker, and the weakness is in what the post asserts about other people. The Christians in question have ``stopped really believing,'' and for them ``No search ever occurred.'' These are claims about the inner states and histories of unnamed people, and the post gives no case. Its own setup shows only that the reverence came first, and the post itself grants that a later check which could change the conclusion would count.

The post also applies its standard to one side. A genuine judgment of the Bible's literary merit, it says, needs a neutral reading of candidate books, which is ``Easy enough if you're not a Christian.'' Its own rankings, that Tolkien's book is ``a vastly superior literary work'' and \textsc{Harry Potter} superior ``as moral philosophy,'' are given without any reading shown. By the post's own reasoning, a reader who has already rejected the Bible has a bottom line too.

In short: a sound argument about reasons found after the fact, with a usable test in the laptop example, applied in the Bible example to unnamed people's inner histories without evidence and to one side of the question only.
\end{afterword}
